\documentclass[UTF8, 12pt, fontset=fandol, oneside]{ctexart}
\usepackage{researchnotes}

\graphicspath{{math/概率论/}}

\title{第 11 章\quad 马尔可夫链}
\author{}
\date{}

\begin{document}

\maketitle

马尔可夫链最早由 Andrey Markov 于 1906 年引入，目的是说明大数定律也可以适用于并不独立的随机变量。为了理解马尔可夫模型从何而来，先考虑一列独立同分布随机变量
\[
  X_0,X_1,\ldots,X_n,\ldots,
\]
其中我们把 $n$ 看作时间。这是第 10 章一直工作的设定；但在为真实世界现象建模时，独立性可能是过于严格的假设，因为它意味着各个 $X_n$ 之间完全不提供信息。另一个极端是允许 $X_n$ 之间有任意相互作用，但这样会使得即便是基本计算也变得很困难。马尔可夫链是一列随机变量，它具有一步依赖性；这个精确含义稍后会定义。因此，马尔可夫链位于完全独立与完全依赖之间，是一种很有用的折中。

自从被发明以来，马尔可夫链已经在大量领域中变得极其重要，例如生物学、博弈论、金融、机器学习和统计物理。它们也广泛用于复杂分布的模拟，相应算法称为\textbf{马尔可夫链蒙特卡罗}（Markov chain Monte Carlo, MCMC）。本章将介绍马尔可夫链及其性质，下一章会考察一些 MCMC 技术的例子。

\section*{11.1\quad 马尔可夫性质与转移矩阵}

马尔可夫链同时“生活”在空间与时间中：$X_n$ 的所有可能取值所成的集合称为\textbf{状态空间}，下标 $n$ 表示过程随时间的演化。马尔可夫链的状态空间可以是离散的，也可以是连续的；时间也可以是离散的或连续的（在连续时间情形中，可以想象对所有实数 $t\ge 0$ 定义过程 $X_t$）。本章只关注\textbf{离散状态、离散时间}的马尔可夫链，并且状态空间是有限的。具体地，我们假设 $X_n$ 的取值在一个有限集合中，通常取为 $\{1,2,\ldots,M\}$ 或 $\{0,1,\ldots,M\}$。

\noindent\textbf{定义 11.1.1（马尔可夫链）}\quad 随机变量序列 $X_0,X_1,X_2,\ldots$ 取值于状态空间 $\{1,2,\ldots,M\}$。若对所有 $n\ge0$，
\[
  P(X_{n+1}=j\mid X_n=i,X_{n-1}=i_{n-1},\ldots,X_0=i_0)
  =
  P(X_{n+1}=j\mid X_n=i),
\]
则称它为一个\textbf{马尔可夫链}。

量 $P(X_{n+1}=j\mid X_n=i)$ 称为从状态 $i$ 到状态 $j$ 的\textbf{转移概率}。在本书中，说到马尔可夫链时，我们会默认它是\textbf{时间齐次}的；这意味着转移概率 $P(X_{n+1}=j\mid X_n=i)$ 对所有时间 $n$ 都相同。不过需要注意，文献中对于应称“时间齐次马尔可夫链”还是直接称“马尔可夫链”并不完全一致。

上面的条件称为\textbf{马尔可夫性质}。它说的是，在给定整个过去历史 $X_0,X_1,X_2,\ldots,X_n$ 的条件下，预测 $X_{n+1}$ 时只有最近的一项 $X_n$ 重要。若把时间 $n$ 看作现在，$n$ 以前的时间看作过去，$n$ 以后的时间看作未来，那么马尔可夫性质就是说：给定现在，过去与未来条件独立。马尔可夫性质极大地简化了条件概率的计算：我们不必以整个过去为条件，只需以最近的取值为条件。

为了描述一个马尔可夫链的动态，我们需要知道从任一状态转移到任一其他状态的概率，即马尔可夫性质右端的概率 $P(X_{n+1}=j\mid X_n=i)$。这些信息可以编码在一个矩阵中，称为\textbf{转移矩阵}；它的 $(i,j)$ 元是链从状态 $i$ 一步转移到状态 $j$ 的概率。

\noindent\textbf{定义 11.1.2（转移矩阵）}\quad 设 $X_0,X_1,X_2,\ldots$ 是状态空间为 $\{1,2,\ldots,M\}$ 的马尔可夫链，并令
\[
  q_{ij}=P(X_{n+1}=j\mid X_n=i)
\]
为从状态 $i$ 到状态 $j$ 的转移概率。$M\times M$ 矩阵 $Q=(q_{ij})$ 称为该链的\textbf{转移矩阵}。

注意，$Q$ 是一个非负矩阵，且每一行的和都是 $1$。这是因为，从任一状态 $i$ 出发，“转移到 1”、“转移到 2”、$\ldots$、“转移到 $M$”这些事件互不相交；而它们的并集概率为 $1$，因为链总要转移到某个地方。

\noindent\textbf{例 11.1.3（雨天--晴天马尔可夫链）}\quad 假设任一天的天气要么是雨天，要么是晴天。若今天下雨，则明天下雨的概率为 $1/3$，晴天的概率为 $2/3$；若今天晴天，则明天下雨的概率为 $1/2$，晴天的概率为 $1/2$。令 $X_n$ 表示第 $n$ 天的天气，则 $X_0,X_1,X_2,\ldots$ 是状态空间 $\{R,S\}$ 上的马尔可夫链，其中 $R$ 表示雨天，$S$ 表示晴天。马尔可夫性质成立，是因为根据过程的描述，只有今天的天气会影响对明天天气的预测。

该链的转移矩阵为
\[
\begin{array}{c@{\;}c}
 & \begin{array}{cc} R & S \end{array}\\[-2mm]
\begin{array}{c} R\\ S \end{array}
&
\begin{pmatrix}
  1/3 & 2/3\\
  1/2 & 1/2
\end{pmatrix}.
\end{array}
\]
第一行表示，从状态 $R$ 出发，以概率 $1/3$ 回到状态 $R$，以概率 $2/3$ 转移到状态 $S$。第二行表示，从状态 $S$ 出发，以概率 $1/2$ 转移到状态 $R$，以概率 $1/2$ 留在状态 $S$。同样也可以把转移矩阵写成
\[
\begin{array}{c@{\;}c}
 & \begin{array}{cc} S & R \end{array}\\[-2mm]
\begin{array}{c} S\\ R \end{array}
&
\begin{pmatrix}
  1/2 & 1/2\\
  2/3 & 1/3
\end{pmatrix}
\end{array}
\]
来表示。在一般情形下，如果马尔可夫链的状态没有显然的顺序（如这里的 $R$ 和 $S$），只需固定一个状态顺序并始终一致地使用它。

马尔可夫链的转移概率也可以用图表示。每个状态用一个圆表示，箭头表示可能的一步转移；可以想象一个粒子在各状态之间游走，并随机选择要跟随的箭头。箭头旁边写出相应的转移概率。

\begin{figure}[htbp]
  \centering
  \includegraphics[width=0.42\textwidth]{11-0-1.png}
  \caption{雨天--晴天马尔可夫链的转移图。}
\end{figure}

如果明天的天气不仅依赖于今天，还依赖于昨天，会怎样？例如，假设天气行为与上面相同，只是若已经连续两天下雨，则明天一定晴天；若已经连续两天晴天，则明天一定下雨。在新的天气动态下，$X_n$ 不再构成马尔可夫链，因为马尔可夫性质被破坏：即使给定今天的天气，昨天的天气仍然能为预测明天提供有用信息。

不过，通过扩大状态空间，我们可以构造一个新的马尔可夫链：令
\[
  Y_n=(X_{n-1},X_n),\qquad n\ge1.
\]
则 $Y_1,Y_2,\ldots$ 是状态空间
\[
  \{(R,R),(R,S),(S,R),(S,S)\}
\]
上的马尔可夫链。可以验证，新的转移矩阵为
\[
\begin{array}{c@{\;}c}
 & \begin{array}{cccc} (R,R) & (R,S) & (S,R) & (S,S) \end{array}\\[-2mm]
\begin{array}{c}
(R,R)\\
(R,S)\\
(S,R)\\
(S,S)
\end{array}
&
\begin{pmatrix}
0 & 1 & 0 & 0\\
0 & 0 & 1/2 & 1/2\\
1/3 & 2/3 & 0 & 0\\
0 & 0 & 1 & 0
\end{pmatrix}
\end{array}
\]
并且其对应的图形表示如下。

\begin{figure}[htbp]
  \centering
  \includegraphics[width=0.36\textwidth]{11-0-2.png}
  \caption{扩充状态空间后的天气马尔可夫链。}
\end{figure}

类似地，通过进一步扩大状态空间，也可以处理天气中的三阶或四阶依赖关系，使马尔可夫性质成立。$\square$

一旦有了马尔可夫链的转移矩阵 $Q$，就可以求出更长时间尺度上的转移概率。

\noindent\textbf{定义 11.1.4（$n$ 步转移概率）}\quad 从 $i$ 到 $j$ 的 \textbf{$n$ 步转移概率}是指，从 $i$ 出发恰好 $n$ 步后处于 $j$ 的概率。记作 $q_{ij}^{(n)}$：
\[
  q_{ij}^{(n)}=P(X_n=j\mid X_0=i).
\]

注意
\[
  q_{ij}^{(2)}=\sum_k q_{ik}q_{kj},
\]
因为要在两步内从 $i$ 到 $j$，链必须先从 $i$ 到某个中间状态 $k$，再从 $k$ 到 $j$；由于马尔可夫性质，这些转移可按上述方式相乘。右端正是矩阵乘法定义下 $Q^2$ 的 $(i,j)$ 元，所以 $Q^2$ 给出两步转移概率。归纳可知，转移矩阵的第 $n$ 次幂给出 $n$ 步转移概率：
\[
  q_{ij}^{(n)} \text{ 是 } Q^n \text{ 的 } (i,j) \text{ 元}.
\]

\noindent\textbf{例 11.1.5（4 状态马尔可夫链的转移矩阵）}\quad 考虑图 11.1 所示的 4 状态马尔可夫链。当箭头上没有写出概率时，如本例中这样，意思是从同一状态出发的所有箭头等可能。例如，从状态 1 出发有 3 条箭头，所以转移 $1\to3$、$1\to2$ 与 $1\to1$ 的概率都是 $1/3$。因此该链的转移矩阵为
\[
  Q=
  \begin{pmatrix}
    1/3 & 1/3 & 1/3 & 0\\
    0 & 0 & 1/2 & 1/2\\
    0 & 1 & 0 & 0\\
    1/2 & 0 & 0 & 1/2
  \end{pmatrix}.
\]

\begin{figure}[htbp]
  \centering
  \includegraphics[width=0.35\textwidth]{11-1.png}
  \caption{一个 4 状态马尔可夫链。}
\end{figure}

为了计算从状态 1 出发，5 步之后链处于状态 3 的概率，应查看 $Q^5$ 的 $(1,3)$ 元。这里用计算机求得
\[
  Q^5=
  \begin{pmatrix}
    853/3888 & 509/1944 & 52/243 & 395/1296\\
    173/864 & 85/432 & 31/108 & 91/288\\
    37/144 & 29/72 & 1/9 & 11/48\\
    499/2592 & 395/1296 & 71/324 & 245/864
  \end{pmatrix},
\]
所以 $q_{13}^{(5)}=52/243$。$\square$

用第 7 章的语言说，转移矩阵 $Q$ 编码了给定链的初始状态时 $X_1$ 的条件分布。具体地，$Q$ 的第 $i$ 行就是给定 $X_0=i$ 时 $X_1$ 的条件 PMF，并以行向量显示。类似地，$Q^n$ 的第 $i$ 行是给定 $X_0=i$ 时 $X_n$ 的条件 PMF。

为了得到 $X_0,X_1,\ldots$ 的边缘分布，除了转移矩阵外，还必须指定链的\textbf{初始条件}。初始状态 $X_0$ 可以是确定的，也可以按某个分布随机给出。令 $(t_1,t_2,\ldots,t_M)$ 为 $X_0$ 的 PMF，并把它写成向量；也就是说，$t_i=P(X_0=i)$。那么链在任意时刻的边缘分布都可以由转移矩阵算出，本质上是用全概率公式对所有状态取平均。

\noindent\textbf{命题 11.1.6（$X_n$ 的边缘分布）}\quad 定义 $\mathbf t=(t_1,t_2,\ldots,t_M)$，其中 $t_i=P(X_0=i)$，并把 $\mathbf t$ 看作行向量。则 $X_n$ 的边缘分布由向量 $\mathbf t Q^n$ 给出。也就是说，$\mathbf t Q^n$ 的第 $j$ 个分量为 $P(X_n=j)$。

\noindent\textit{证明.}\quad 由全概率公式，并对 $X_0$ 条件化，链在 $n$ 步后处于状态 $j$ 的概率为
\[
  P(X_n=j)
  =
  \sum_{i=1}^M P(X_0=i)P(X_n=j\mid X_0=i)
  =
  \sum_{i=1}^M t_i q_{ij}^{(n)},
\]
这正是矩阵乘法定义下 $\mathbf tQ^n$ 的第 $j$ 个分量。$\blacksquare$

\noindent\textbf{例 11.1.7（4 状态马尔可夫链的边缘分布）}\quad 再次考虑图 11.1 中的 4 状态马尔可夫链。假设初始条件为
\[
  \mathbf t=(1/4,1/4,1/4,1/4),
\]
即链以相同概率从四个状态中的任一个开始。令 $X_n$ 表示链在时刻 $n$ 的位置。则 $X_1$ 的边缘分布为
\[
\begin{aligned}
  \mathbf tQ
  &=
  \begin{pmatrix}1/4&1/4&1/4&1/4\end{pmatrix}
  \begin{pmatrix}
    1/3 & 1/3 & 1/3 & 0\\
    0 & 0 & 1/2 & 1/2\\
    0 & 1 & 0 & 0\\
    1/2 & 0 & 0 & 1/2
  \end{pmatrix}\\
  &=
  \begin{pmatrix}5/24&1/3&5/24&1/4\end{pmatrix}.
\end{aligned}
\]
$X_5$ 的边缘分布为
\[
\begin{aligned}
  \mathbf tQ^5
  &=
  \begin{pmatrix}1/4&1/4&1/4&1/4\end{pmatrix}
  \begin{pmatrix}
    853/3888 & 509/1944 & 52/243 & 395/1296\\
    173/864 & 85/432 & 31/108 & 91/288\\
    37/144 & 29/72 & 1/9 & 11/48\\
    499/2592 & 395/1296 & 71/324 & 245/864
  \end{pmatrix}\\
  &=
  \begin{pmatrix}3379/15552&2267/7776&101/486&1469/5184\end{pmatrix}.
\end{aligned}
\]
这里我们用计算机完成了矩阵乘法。$\square$

\section*{11.2\quad 状态分类}

本节引入一些术语，用来描述马尔可夫链的各种特征。马尔可夫链中的状态可以分为\textbf{常返}或\textbf{暂留}，这取决于长期来看它们会被反复访问，还是最终被抛弃。状态也可以按其\textbf{周期}分类；周期是一个正整数，用来概括连续两次访问同一状态之间可能经过的时间量。这些特征很重要，因为它们决定了马尔可夫链的长期行为，这将在第 11.3 节研究。

常返和暂留的概念最好通过具体例子来说明。在图 11.2 左侧的马尔可夫链中（它先前出现在例 11.1.5 中），一个在状态之间移动的粒子长期来看会继续在所有 4 个状态中停留，因为从任一状态都可能到达任一其他状态。相比之下，考虑图 11.2 右侧的链，并令粒子从状态 1 开始。链可能会在由状态 1、2、3 形成的三角形中逗留一段时间，但最终会到达状态 4；从那以后，它再也不能回到状态 1、2 或 3，而会永远在状态 4、5、6 之间游走。因此，状态 1、2、3 是暂留的，状态 4、5、6 是常返的。

\begin{figure}[htbp]
  \centering
  \includegraphics[width=0.58\textwidth]{11-2.png}
  \caption{左：所有状态都是常返的 4 状态马尔可夫链。右：状态 1、2、3 为暂留状态的 6 状态马尔可夫链。}
\end{figure}

一般地，这些概念定义如下。

\noindent\textbf{定义 11.2.1（常返状态与暂留状态）}\quad 马尔可夫链的状态 $i$ 称为\textbf{常返}的，如果从 $i$ 出发，链最终返回 $i$ 的概率为 $1$。否则称该状态为\textbf{暂留}的；这意味着若链从 $i$ 出发，则有正概率永远不返回 $i$。

事实上，虽然暂留状态的定义只要求“有正概率永远不返回该状态”，但可以说得更强：只要有正概率永远离开 $i$，链最终就一定会永远离开 $i$。此外，我们还可以求出返回该状态次数的分布。

\noindent\textbf{命题 11.2.2（返回暂留状态的次数服从几何分布）}\quad 令 $i$ 为马尔可夫链的一个暂留状态。假设从 $i$ 出发永不返回 $i$ 的概率为正数 $p>0$。则从 $i$ 出发，在永远离开之前链返回 $i$ 的次数服从 $\operatorname{Geom}(p)$ 分布。

证明来自几何分布的故事：每次链处于 $i$ 时，都有一次 Bernoulli 试验；若链最终返回 $i$，记为“失败”，若链永远离开 $i$，记为“成功”。这些试验由于马尔可夫性质而独立。返回状态 $i$ 的次数就是第一次成功之前失败的次数，正好符合几何分布的故事。又因为几何随机变量总是取有限值，该命题告诉我们，经过有限次访问后，链将永远离开状态 $i$。

当状态数不太大时，分类状态为常返或暂留的一种方法是画出马尔可夫链的图，并使用分析图 11.2 中链时同样的推理。一个可以立刻断定所有状态均为常返的特殊情形，是链为\textbf{不可约}，意思是从任一状态都可以到达任一其他状态。

\noindent\textbf{定义 11.2.3（不可约链与可约链）}\quad 转移矩阵为 $Q$ 的马尔可夫链称为\textbf{不可约}的，如果对任意两个状态 $i$ 和 $j$，都可以在有限步内以正概率从 $i$ 到达 $j$。也就是说，对任意状态 $i,j$，存在正整数 $n$，使得 $Q^n$ 的 $(i,j)$ 元为正。不是不可约的马尔可夫链称为\textbf{可约}的。

\noindent\textbf{命题 11.2.4（不可约推出所有状态常返）}\quad 在有限状态空间的不可约马尔可夫链中，所有状态都是常返的。

\noindent\textit{证明.}\quad 显然至少有一个状态必须是常返的；如果所有状态都是暂留的，链最终会永远离开所有状态，从而无处可去。于是无妨假设状态 1 是常返的，并考虑任一其他状态 $i$。由不可约的定义，存在某个 $n$ 使得 $q_{1i}^{(n)}$ 为正。因此，链每次处于状态 1 时，都有正概率在再经过 $n$ 步后到达状态 $i$。

由于链无限多次访问状态 1，我们知道它最终一定会从状态 1 到达状态 $i$；可以把每次访问状态 1 看作一次试验，其中“成功”定义为在至多 $n$ 步内到达状态 $i$。从状态 $i$ 出发，链会返回状态 1，因为状态 1 是常返的；按同样逻辑，它之后又会再次到达状态 $i$。归纳可知，链将无限多次访问状态 $i$。由于 $i$ 是任意的，所有状态都是常返的。$\blacksquare$

该命题的逆命题是假的：一个可约马尔可夫链也可能所有状态都是常返的。下面这个马尔可夫链给出了一个例子，它由两个状态“岛”组成。

\begin{figure}[htbp]
  \centering
  \includegraphics[width=0.48\textwidth]{11-2-1.png}
  \caption{由两个常返类组成的可约马尔可夫链。}
\end{figure}

\noindent\textbf{注意 11.2.5}\quad 注意，常返或暂留是马尔可夫链中每个\textbf{状态}的性质，而不可约或可约是整条链的性质。

下面从马尔可夫链的角度重新看两个前几章熟悉的问题。对每一个问题，我们会辨认常返状态和暂留状态。

\noindent\textbf{例 11.2.6（作为马尔可夫链的赌徒破产）}\quad 在赌徒破产问题（例 2.7.3）中，两个赌徒 A 与 B 分别以 $i$ 美元和 $N-i$ 美元开始，并连续进行每局 1 美元的下注。每轮中，玩家 A 以概率 $p$ 赢，以概率 $q=1-p$ 输。令 $X_n$ 为赌徒 A 在时刻 $n$ 的财富。则 $X_0,X_1,\ldots$ 是状态空间 $\{0,1,\ldots,N\}$ 上的马尔可夫链。按设定，$X_0=i$。一旦马尔可夫链到达 $0$ 或 $N$，分别表示玩家 A 或玩家 B 破产，链就会永远停留在该状态。该链的图如下。

\begin{figure}[htbp]
  \centering
  \includegraphics[width=0.62\textwidth]{11-2-2.png}
  \caption{赌徒破产问题对应的马尔可夫链。}
\end{figure}

我们在第 2 章已经证明，A 或 B 中必有一人破产的概率为 $1$。因此，对任意不是 $0$ 或 $N$ 的初始状态 $i$，该马尔可夫链最终会被吸收到状态 $0$ 或 $N$，而不再返回 $i$。所以，对这条马尔可夫链，状态 $0$ 和 $N$ 是常返的，所有其他状态都是暂留的。该链是可约的，因为从状态 $0$ 只能到状态 $0$，从状态 $N$ 只能到状态 $N$。$\square$

\noindent\textbf{例 11.2.7（作为马尔可夫链的收集优惠券问题）}\quad 在收集优惠券问题（例 4.3.12）中，共有 $C$ 种优惠券（或玩具），我们一次收集一种，每次都从 $C$ 种优惠券中有放回抽样。令 $X_n$ 为 $n$ 次尝试后，我们收藏中不同优惠券类型的数量。则 $X_0,X_1,\ldots$ 是状态空间 $\{0,1,\ldots,C\}$ 上的马尔可夫链。按设定，$X_0=0$。该链如下图所示。

\begin{figure}[htbp]
  \centering
  \includegraphics[width=0.62\textwidth]{11-2-3.png}
  \caption{收集优惠券问题对应的马尔可夫链。}
\end{figure}

除了状态 $C$ 之外，一旦离开某个状态，我们就永远不能返回；收藏中的优惠券种类数只会随时间增加。因此，除 $C$ 外所有状态都是暂留的，而 $C$ 是常返的。该链是可约的，例如不可能从状态 2 回到状态 1。$\square$

另一种分类状态的方法，是按照状态的周期。一个状态的周期概括了连续两次访问该状态之间可能经过多少时间。

\noindent\textbf{定义 11.2.8（状态的周期，周期链与非周期链）}\quad 马尔可夫链中状态 $i$ 的\textbf{周期}是指，从 $i$ 出发返回 $i$ 时可能经历的步数集合的最大公约数。也就是说，$i$ 的周期是所有满足 $Q^n$ 的 $(i,i)$ 元为正的 $n$ 的最大公约数。（若从 $i$ 出发永远不可能返回 $i$，则 $i$ 的周期为 $\infty$。）若某状态的周期为 $1$，则称它是\textbf{非周期}的；否则称为\textbf{周期}的。若链的所有状态都是非周期的，则称这条链是\textbf{非周期}的；否则称为\textbf{周期}的。

例如，再考虑图 11.2 中的两条马尔可夫链，它们在图 11.3 中再次画出。先看右侧的 6 状态链。从状态 1 出发，可能在 3 步、6 步、9 步等之后回到状态 1，但不可能在不是 3 的倍数的步数后回到状态 1。因此，状态 1 的周期为 3。类似地，状态 2 和 3 的周期也为 3。另一方面，状态 4、5、6 的周期为 1；但由于至少有一个状态的周期不是 1，该链是周期的。相比之下，左侧链的所有状态都是非周期的，所以它是非周期链。

\begin{figure}[htbp]
  \centering
  \includegraphics[width=0.58\textwidth]{11-3.png}
  \caption{左：一个非周期马尔可夫链。右：一个周期马尔可夫链，其中状态 1、2、3 的周期为 3。}
\end{figure}

在例 11.2.6 的赌徒破产链中，除状态 $0$ 和 $N$ 外，每个状态的周期都是 2，而状态 $0$ 与 $N$ 的周期为 1。在收集优惠券链中，除状态 0 外每个状态的周期都是 1；状态 0 的周期为 $\infty$，因为不可能返回状态 0。因此这两条链都不是非周期链。

检查不可约链是否为非周期，常常比想象中容易得多：下一个命题说明，只需要计算一个状态的周期，而不必逐个状态寻找周期不是 1 的状态。

\noindent\textbf{命题 11.2.9（不可约链中的周期）}\quad 在不可约马尔可夫链中，所有状态有相同周期。

\section*{11.3\quad 平稳分布}

常返与暂留的概念对于理解马尔可夫链的长期行为很重要。起初，链可能会在暂留状态中花时间；但最终，链会把全部时间花在常返状态上。那么，在每个常返状态中，它会花多大比例的时间？这个问题由链的\textbf{平稳分布}回答，也称为\textbf{稳态分布}。本节会看到，对于不可约且非周期的马尔可夫链，平稳分布描述了链的长期行为，而与初始条件无关。它会告诉我们两件事：长期来看处于任一特定状态的概率，以及链在该状态中花费时间的长期比例。

\noindent\textbf{定义 11.3.1（平稳分布）}\quad 设 $\mathbf s=(s_1,\ldots,s_M)$ 为行向量，满足 $s_i\ge0$ 且 $\sum_i s_i=1$。若对转移矩阵为 $Q$ 的马尔可夫链，有
\[
  \sum_i s_i q_{ij}=s_j
\]
对所有 $j$ 成立，则称 $\mathbf s$ 是该链的\textbf{平稳分布}。这个线性方程组可以写成一个矩阵方程：
\[
  \mathbf sQ=\mathbf s.
\]

回忆一下，如果 $\mathbf s$ 是 $X_0$ 的分布，那么 $\mathbf sQ$ 是 $X_1$ 的边缘分布。因此方程 $\mathbf sQ=\mathbf s$ 的含义是：如果 $X_0$ 的分布为 $\mathbf s$，那么 $X_1$ 的分布也为 $\mathbf s$；于是 $X_2$、$X_3$ 等也都有分布 $\mathbf s$。也就是说，若马尔可夫链的初始分布就是平稳分布 $\mathbf s$，它就会永远保持在平稳分布中。

直观理解马尔可夫链的平稳分布的一种方式，是想象有大量粒子，它们各自独立地按照转移概率在状态之间跳动。经过一段时间后，这个粒子系统会趋近于一种平衡：在每一个时间段内，离开某状态的粒子数会被进入该状态的粒子数抵消，而且所有状态都如此。在这种平衡下，整个系统看起来是平稳的，而每个状态中的粒子比例就由平稳分布给出。在定义 11.4.1 之后，我们会进一步探索这种关于平稳分布的观点。

\noindent\textbf{注意 11.3.2（平稳分布是边缘的，不是条件的）}\quad 当马尔可夫链处于平稳分布时，$X_n$ 的无条件 PMF 对所有 $n$ 都等于 $\mathbf s$；但给定 $X_{n-1}=i$ 时 $X_n$ 的条件 PMF 仍由转移矩阵 $Q$ 的第 $i$ 行编码。

如果一个马尔可夫链从平稳分布开始，那么所有 $X_n$ 都同分布（因为它们有相同的边缘分布 $\mathbf s$），但它们未必独立，因为给定 $X_{n-1}=i$ 时 $X_n$ 的条件分布一般不同于 $X_n$ 的边缘分布。

\noindent\textbf{注意 11.3.3（同情魔法）}\quad 如果一个马尔可夫链从平稳分布开始，那么 $X_n$ 的边缘分布都相等。这并不等于说随机变量 $X_n$ 本身都相等；把随机变量与其分布混淆，就是一种“同情魔法”的例子。

对于很小的马尔可夫链，可以直接用定义手算平稳分布。下一个例子用一个两状态链说明这一点。

\noindent\textbf{例 11.3.4（两状态链的平稳分布）}\quad 令
\[
  Q=
  \begin{pmatrix}
    1/3 & 2/3\\
    1/2 & 1/2
  \end{pmatrix}.
\]
平稳分布形如 $\mathbf s=(s,1-s)$，于是需要解
\[
  \begin{pmatrix}s&1-s\end{pmatrix}
  \begin{pmatrix}
    1/3 & 2/3\\
    1/2 & 1/2
  \end{pmatrix}
  =
  \begin{pmatrix}s&1-s\end{pmatrix},
\]
这等价于
\[
  \frac13s+\frac12(1-s)=s,
\]
\[
  \frac23s+\frac12(1-s)=1-s.
\]
这些方程的唯一解是 $s=3/7$，所以 $(3/7,4/7)$ 是该马尔可夫链唯一的平稳分布。

更一般地，假设 $q_{12}=a$ 且 $q_{21}=b$，其中 $0<a<1$ 且 $0<b<1$。则转移矩阵为
\[
  Q=
  \begin{pmatrix}
    1-a & a\\
    b & 1-b
  \end{pmatrix}.
\]
写 $\mathbf s=(s_1,s_2)$，方程 $\mathbf sQ=\mathbf s$ 变成线性方程组
\[
  (1-a)s_1+bs_2=s_1,
\]
\[
  as_1+(1-b)s_2=s_2.
\]
这两个方程都化简为
\[
  as_1=bs_2.
\]
代入 $s_2=1-s_1$，可得该系统的唯一解为
\[
  \mathbf s=\left(\frac{b}{a+b},\frac{a}{a+b}\right).
\]
简言之，$\mathbf s\propto(b,a)$，也就是说 $\mathbf s$ 是 $(b,a)$ 的某个常数倍；这个常数由“分量和为 1”的要求决定。作为检验，对前面给出的具体 $Q$，这个结果说明平稳分布应与 $(b,a)=(1/2,2/3)$ 成比例；乘以 6 清除分母后，相当于与 $(3,4)$ 成比例。前面求出的平稳分布为 $(3/7,4/7)$，确实与 $(3,4)$ 成比例。$\square$

在线性代数术语中，方程 $\mathbf sQ=\mathbf s$ 表示 $\mathbf s$ 是 $Q$ 的左特征向量，对应特征值 1（参见数学附录 A.3）。为了得到通常意义下的特征向量（右特征向量），可以取转置：$Q'\mathbf s'=\mathbf s'$，其中撇号表示转置。

\subsection*{11.3.1\quad 存在唯一性}

平稳分布是否总是存在？是否唯一？事实证明，对有限状态空间，平稳分布总是存在。此外，在不可约马尔可夫链中，平稳分布唯一。

\noindent\textbf{定理 11.3.5（平稳分布的存在唯一性）}\quad 对任意不可约马尔可夫链，存在唯一的平稳分布。在该分布中，每个状态都有正概率。

这个定理是线性代数中 Perron--Frobenius 定理的一个推论；该定理陈述于数学附录 A.3。

图 11.3 左侧的 4 状态链是不可约的：从图形角度看，沿着箭头可以从任何地方到达任何地方；从转移矩阵 $Q$ 的角度看，$Q^5$ 的所有元素都为正。因此，由定理 11.3.5，该链有唯一平稳分布。

另一方面，例 11.2.6 中的赌徒破产链是可约的，所以该定理不适用。事实证明，这条链有无穷多个平稳分布。长期来看，链最终会停在状态 0（并永远留在那里），或停在状态 $N$（并永远留在那里）。这提示我们，而且也很容易检验，退化分布
\[
  \mathbf s=(1,0,\ldots,0),\qquad \mathbf t=(0,0,\ldots,1)
\]
都是平稳分布。因此，对任意 $0\le p\le1$，任意加权组合 $p\mathbf s+(1-p)\mathbf t$ 也是平稳分布，因为它的分量和为 1，且
\[
  (p\mathbf s+(1-p)\mathbf t)Q
  =
  p\mathbf sQ+(1-p)\mathbf tQ
  =
  p\mathbf s+(1-p)\mathbf t,
\]
其中照例 $Q$ 是转移矩阵。

\subsection*{11.3.2\quad 收敛性}

我们已经非正式地说过，平稳分布描述了链的长期行为，其含义是，如果让链运行很久，$X_n$ 的边缘分布会收敛到平稳分布 $\mathbf s$。下一个定理说明，只要链既不可约又非周期，这一点就成立。于是，无论链的初始条件如何，随着 $n\to\infty$，$X_n$ 的 PMF 都会收敛到平稳分布。这把平稳性的概念与马尔可夫链的长期行为联系起来。证明略去。

\noindent\textbf{定理 11.3.6（收敛到平稳分布）}\quad 设 $X_0,X_1,\ldots$ 是不可约、非周期马尔可夫链，平稳分布为 $\mathbf s$，转移矩阵为 $Q$。则当 $n\to\infty$ 时，$P(X_n=i)$ 收敛到 $s_i$。用转移矩阵表示，$Q^n$ 收敛到一个每一行都是 $\mathbf s$ 的矩阵。

因此，经过大量步之后，链处于状态 $i$ 的概率接近平稳概率 $s_i$，而与链的初始条件无关。这使得不可约、非周期链特别好处理。不可约意味着对每个 $(i,j)$，存在某个幂次 $Q^m$ 使其 $(i,j)$ 元为正；如果再假设非周期，事实证明可以找到一个对所有 $i,j$ 同时适用的 $m$。更精确地说，一条链不可约且非周期，当且仅当存在某个幂次 $Q^m$ 的所有元素都为正。

直观地说，额外的非周期条件是为了排除那些只会绕圈运行的链，例如下面例子中的链；也为了排除那些只能在偶数步后到达某些状态、只能在奇数步后到达另一些状态的链。

\noindent\textbf{例 11.3.7（周期链）}\quad 下图给出了一条周期马尔可夫链，其中每个状态的周期都是 5。

\begin{figure}[htbp]
  \centering
  \includegraphics[width=0.28\textwidth]{11-4.png}
  \caption{一条周期链。}
\end{figure}

该链的转移矩阵为
\[
  Q=
  \begin{pmatrix}
    0 & 1 & 0 & 0 & 0\\
    0 & 0 & 1 & 0 & 0\\
    0 & 0 & 0 & 1 & 0\\
    0 & 0 & 0 & 0 & 1\\
    1 & 0 & 0 & 0 & 0
  \end{pmatrix}.
\]
不难验证，$\mathbf s=(1/5,1/5,1/5,1/5,1/5)$ 是这条链的平稳分布；由定理 11.3.5，它是唯一的。不过，假设链从 $X_0=1$ 开始，则 $X_n$ 的 PMF 对状态 $(n\bmod 5)+1$ 赋予概率 1，对其他所有状态赋予概率 0。\footnote{记号 $n\bmod 5$ 表示 $n$ 除以 5 后的余数。}特别地，当 $n\to\infty$ 时，它并不收敛到 $\mathbf s$。$Q^n$ 也不会收敛到一个每一行都是 $\mathbf s$ 的矩阵：该链的转移是确定性的，所以 $Q^n$ 总是由 0 和 1 组成。$\square$

最后，平稳分布还告诉我们访问任一特定状态之间的平均时间。

\noindent\textbf{定理 11.3.8（期望返回时间）}\quad 设 $X_0,X_1,\ldots$ 是不可约马尔可夫链，平稳分布为 $\mathbf s$。令 $r_i$ 表示在链从 $i$ 出发的条件下，链返回 $i$ 所需时间的期望。则
\[
  s_i=\frac{1}{r_i}.
\]

下面说明这些定理如何应用到例 11.3.4 中的两状态链。

\noindent\textbf{例 11.3.9（两状态链的长期行为）}\quad 长期来看，例 11.3.4 中的链会把 $3/7$ 的时间花在状态 1，把 $4/7$ 的时间花在状态 2。从状态 1 开始，平均需要 $7/3$ 步返回状态 1。转移矩阵的幂收敛到一个每一行都是平稳分布的矩阵：
\[
  \begin{pmatrix}
    1/3 & 2/3\\
    1/2 & 1/2
  \end{pmatrix}^n
  \to
  \begin{pmatrix}
    3/7 & 4/7\\
    3/7 & 4/7
  \end{pmatrix},
  \qquad n\to\infty.
\]
$\square$

\subsection*{11.3.3\quad Google PageRank}

接下来考虑一个规模大得多的平稳分布例子：状态空间包含数十亿个相互连接节点的马尔可夫链，也就是万维网。下一个例子解释了 Google 的创始人如何把网上冲浪建模为马尔可夫链，并用其平稳分布来排序网页相关性。多年里，Google 把由此得到的方法称为 PageRank，并称它为“我们软件的核心”。

假设你对某个主题感兴趣，比如国际象棋，于是使用搜索引擎查找含有国际象棋信息的有用网页。有数百万个网页提到“chess”这个词，所以搜索引擎需要处理的关键问题是：搜索结果应以什么顺序显示？如果必须先翻过成千上万个仅仅提到“chess”的垃圾页面，才能找到有信息量的内容，那将是一场灾难。

在网络早期，人们尝试过各种排序方法。例如，有些搜索引擎雇人手工判断哪些页面最有用，像博物馆策展人一样。但除了主观且昂贵之外，随着网络增长，这很快变得不可行。另一些方法关注搜索词在网站上出现的次数。但一个反复提到“chess”的页面，可能很容易不如一个简明的参考页面，或一个关于国际象棋但没有反复提到这个词的页面有用。此外，这种方法极易被滥用：垃圾页面只要包含很长的重复词表，就能提升自己的排名。

上述两种方法都忽略了网络的\textbf{结构}：哪些页面链接到哪些其他页面？把链接结构纳入考虑，使搜索引擎出现了巨大改进。作为第一次尝试，可以按有多少其他页面链接到某个页面来给该页面排序。也就是说，如果页面 A 链接到页面 B，我们就把它看作给 B 投了一票，并根据页面拥有多少票来排序。

但这同样非常容易被滥用：垃圾页面可以通过创建成千上万个链接到自己的其他垃圾页面来提高排名。并且，虽然让每个页面有相同投票权似乎很民主，但来自可靠页面的入链比来自信息贫乏页面的链接更有意义。Google PageRank 由 Sergey Brin 和名字恰如其分的 Larry Page 于 1998 年提出，它不仅按有多少页面链接到某页面来衡量该页面的重要性，也按这些页面本身的重要性来衡量。

把网络看作一个有向网络；事实上它也正是如此。网络中的每个页面是一个节点，节点之间的链接表示页面之间的链接。比如，为简单起见，假设网络只有 4 个页面，其连接关系如下图所示。

\begin{figure}[htbp]
  \centering
  \includegraphics[width=0.35\textwidth]{11-5.png}
  \caption{一个很小的网络。}
\end{figure}

想象某个人在网上随机冲浪：从某个页面开始，然后随机点击链接从一个页面跳到下一个页面（在当前页面的所有链接之间等可能选择）。PageRank 的想法是：用长期来看停留在某页面上的时间比例来衡量该页面的重要性。

当然，有些页面可能完全没有外链，比如上面的页面 4。网络冲浪者遇到这样的页面时，不会陷入绝望，而是打开一个新的浏览器窗口并访问一个均匀随机页面。因此，一个没有链接的页面会被转换为一个链接到所有页面（包括它自己）的页面。对上面的例子，由此得到的转移矩阵为
\[
  Q=
  \begin{pmatrix}
    0 & 1/2 & 0 & 1/2\\
    1/2 & 0 & 1/2 & 0\\
    0 & 0 & 0 & 1\\
    1/4 & 1/4 & 1/4 & 1/4
  \end{pmatrix}.
\]

一般地，令 $M$ 为网页数量，令 $Q$ 为上述链的 $M\times M$ 转移矩阵，并令 $\mathbf s$ 为平稳分布（假设它存在且唯一）。可以把 $s_j$ 看作页面 $j$ 的重要性度量。直观地，方程
\[
  s_j=\sum_i s_iq_{ij}
\]
说明，页面 $j$ 的分数不仅应基于有多少其他页面链接到它，还应基于这些页面的分数。此外，一个页面的“投票权”会被它的大量外链稀释：如果页面 $i$ 唯一的链接指向页面 $j$（所以 $q_{ij}=1$），这比页面 $i$ 有数千个链接且其中一个恰好指向 $j$ 时更重要。

对于这条链，平稳分布是否唯一并不明显，因为它可能并非不可约且非周期。即使它不可约且非周期，由于网络极其庞大，收敛到平稳分布的速度也可能很慢。为了解决这些问题，假设每次移动前，网络冲浪者以概率 $\alpha$ 掷出正面。如果是正面，冲浪者点击当前页面中的一个随机链接；如果是反面，冲浪者\textbf{传送}到一个均匀随机页面。得到的链具有 Google 转移矩阵
\[
  G=\alpha Q+(1-\alpha)\frac{J}{M},
\]
其中 $J$ 是全 1 的 $M\times M$ 矩阵。注意 $G$ 的每一行和为 1，并且所有元素都为正，所以 $G$ 是一个合法的转移矩阵，并对应一条不可约、非周期的马尔可夫链。这意味着它有唯一平稳分布 $\mathbf s$，称为 PageRank，而且链会收敛到它！$\alpha$ 的选取是一个重要问题；取接近 1 的 $\alpha$ 可以尽可能尊重网络结构，但也有折中，因为较小的 $\alpha$ 会让链收敛得快得多。作为折中，Brin 和 Page 最初推荐 $\alpha=0.85$。

PageRank 在概念上很漂亮，但计算它听起来极其困难，因为 $\mathbf sG=\mathbf s$ 可能是一个含有 1000 亿个方程、1000 亿个未知量的线性系统。与其把它看作一个巨大的代数问题，不如使用马尔可夫链解释：对任意初始分布 $\mathbf t$，有 $\mathbf tG^n\to\mathbf s$，其中 $n\to\infty$。并且 $\mathbf tG$ 起初看起来并没有那么难算：
\[
  \mathbf tG
  =
  \alpha(\mathbf tQ)+\frac{1-\alpha}{M}(\mathbf tJ),
\]
其中第一项并不太难计算，因为 $Q$ 非常稀疏（大多为 0）；第二项也容易，因为 $\mathbf tJ$ 是全 1 向量。然后把 $\mathbf tG$ 作为新的 $\mathbf t$，继续计算 $\mathbf tG^2=(\mathbf tG)G$，如此往复，直到序列看起来已经收敛（虽然很难\textbf{知道}它确实已经收敛）。这给出了 PageRank 的近似值，并且有一个直观解释：它是在大量步之后网络冲浪者所在位置的分布。

\section*{11.4\quad 可逆性}

我们已经看到，马尔可夫链的平稳分布对于理解其长期行为非常有用。遗憾的是，一般来说，当状态空间很大时，求平稳分布在计算上可能很困难。本节讨论一个重要的特殊情形，在这种情形下可以避免处理大矩阵的特征值方程。

\noindent\textbf{定义 11.4.1（可逆性）}\quad 令 $Q=(q_{ij})$ 为马尔可夫链的转移矩阵。假设存在 $\mathbf s=(s_1,\ldots,s_M)$，其中 $s_i\ge0$ 且 $\sum_i s_i=1$，使得
\[
  s_iq_{ij}=s_jq_{ji}
\]
对所有状态 $i$ 与 $j$ 成立。这个方程称为\textbf{可逆性条件}或\textbf{细致平衡条件}；若它成立，则称该链相对于 $\mathbf s$ 是\textbf{可逆}的。

“可逆”一词来自这样一个事实：如果一条可逆链从其平稳分布开始，那么无论时间正向运行还是反向运行，它的行为方式都是一样的。如果你录下一条从平稳分布开始的可逆链的视频，然后把视频放给朋友看，无论按正常顺序播放还是倒放，你的朋友都不能仅凭观看判断时间是在正向还是反向运行。

如定义 11.3.1 之后所讨论的，我们可以直观地把马尔可夫链的平稳分布理解为一个由大量粒子组成的系统；这些粒子按照转移概率独立地在状态间跳动。长期来看，处于任一状态 $j$ 的粒子比例就是状态 $j$ 的平稳概率，而且从状态 $j$ 流出的粒子流会被流入状态 $j$ 的粒子流抵消。更具体地，令 $n$ 为粒子数，令 $\mathbf s$ 为概率向量，使得 $s_j$ 是当前处于状态 $j$ 的粒子比例。按定义，$\mathbf s$ 是该链的平稳分布当且仅当
\[
  s_j=\sum_i s_iq_{ij}=s_jq_{jj}+\sum_{i:i\ne j}s_iq_{ij}
\]
对所有状态 $j$ 成立。这个方程可以改写为
\[
  ns_j(1-q_{jj})=\sum_{i:i\ne j}ns_iq_{ij}.
\]
左端是在下一步将从状态 $j$ 离开的粒子数的近似值，因为状态 $j$ 中有 $ns_j$ 个粒子，每个粒子以概率 $q_{jj}$ 留在 $j$，以概率 $1-q_{jj}$ 离开。右端是在下一步将进入状态 $j$ 的粒子数的近似值，因为对每个 $i\ne j$，状态 $i$ 中有 $ns_i$ 个粒子，每个粒子以概率 $q_{ij}$ 进入状态 $j$。因此，离开状态 $j$ 的粒子与进入状态 $j$ 的粒子之间存在平衡。

可逆性条件施加了一种更强的平衡形式：对每一对状态 $i,j$（$i\ne j$），从状态 $i$ 流向状态 $j$ 的粒子流都被从状态 $j$ 流向状态 $i$ 的粒子流抵消。为看出这一点，把状态 $i$ 与 $j$ 的可逆性方程写作
\[
  ns_iq_{ij}=ns_jq_{ji}.
\]
左端是在下一步将从状态 $i$ 到状态 $j$ 的粒子数的近似值，因为状态 $i$ 中有 $ns_i$ 个粒子，每个粒子以概率 $q_{ij}$ 到状态 $j$。类似地，右端是在下一步将从状态 $j$ 到状态 $i$ 的粒子数的近似值。所以可逆性意味着，从状态 $i$ 到状态 $j$ 的粒子与从状态 $j$ 到状态 $i$ 的粒子之间存在平衡。

给定一个转移矩阵，如果能找到满足可逆性条件的概率向量 $\mathbf s$，那么 $\mathbf s$ 自动就是平稳分布。结合上面的讨论，这并不意外。

\noindent\textbf{命题 11.4.2（可逆推出平稳）}\quad 假设 $Q=(q_{ij})$ 是一条马尔可夫链的转移矩阵，并且该链相对于一个非负向量 $\mathbf s=(s_1,\ldots,s_M)$ 可逆，其中 $\mathbf s$ 的分量和为 1。则 $\mathbf s$ 是该链的平稳分布。

\noindent\textit{证明.}\quad 我们有
\[
  \sum_i s_iq_{ij}=\sum_i s_jq_{ji}=s_j\sum_i q_{ji}=s_j,
\]
其中最后一个等号成立是因为 $Q$ 的每一行和为 1。因此 $\mathbf s$ 是平稳的。$\blacksquare$

这是一个很有力的结果，因为检验可逆性条件往往比解整个方程组 $\mathbf sQ=\mathbf s$ 更容易。不过，一般来说，我们事先未必知道是否可能找到满足可逆性条件的 $\mathbf s$；即使可能，也可能需要很大力气才能找到可用的 $\mathbf s$。在本节余下部分，我们考察三类确实能够找到满足可逆性条件的 $\mathbf s$ 的马尔可夫链。这样的马尔可夫链称为可逆链。

首先，若 $Q$ 是对称矩阵，则平稳分布是在状态空间上的均匀分布：
\[
  \mathbf s=(1/M,1/M,\ldots,1/M).
\]
很容易看出，如果 $q_{ij}=q_{ji}$，那么当所有 $i,j$ 都有 $s_i=s_j$ 时，可逆性条件 $s_iq_{ij}=s_jq_{ji}$ 成立。

这是下面更一般事实的特殊情形：如果 $Q$ 的每一列和为 1，则平稳分布是在状态空间上的均匀分布。

\noindent\textbf{命题 11.4.3}\quad 如果转移矩阵 $Q$ 的每一列和为 1，则所有状态上的均匀分布
\[
  (1/M,1/M,\ldots,1/M)
\]
是一个平稳分布。（行和与列和都为 1 的非负矩阵称为\textbf{双随机矩阵}。）

\noindent\textit{证明.}\quad 假设每一列和为 1，则行向量 $\mathbf v=(1,1,\ldots,1)$ 满足 $\mathbf vQ=\mathbf v$。因此 $(1/M,1/M,\ldots,1/M)$ 是平稳的。$\blacksquare$

其次，如果马尔可夫链是\textbf{无向网络上的随机游走}，那么平稳分布有一个简单公式。

\noindent\textbf{例 11.4.4（无向网络上的随机游走）}\quad 网络是由边连接起来的节点集合；若边可以双向通过，也就是说没有单行道，则称该网络是无向的。假设一个游走者随机沿着无向网络的边移动。从节点 $i$ 出发，游走者等可能地随机选择节点 $i$ 上的任一条边，然后沿选中的边移动。例如，在下图所示网络中，从节点 3 出发，游走者以各 $1/2$ 的概率到节点 1 或节点 2。

\begin{figure}[htbp]
  \centering
  \includegraphics[width=0.32\textwidth]{11-5-1.png}
\end{figure}

一个节点的\textbf{度}是连接到该节点的边数；节点为 $1,2,\ldots,n$ 的网络的\textbf{度序列}是列出所有度的向量 $(d_1,\ldots,d_n)$，其中 $d_j$ 是节点 $j$ 的度。允许从节点到其自身的边（这样的边称为\textbf{自环}），并且它对该节点的度贡献 1。

例如，上面的网络有度序列 $\mathbf d=(4,3,2,3,2)$。注意
\[
  d_iq_{ij}=d_jq_{ji}
\]
对所有 $i,j$ 成立，因为当 $i\ne j$ 时，若 $\{i,j\}$ 是一条边，则 $q_{ij}=1/d_i$，否则为 0。因此，由命题 11.4.2，平稳分布与度序列成比例。直观地，度最高的节点连接最多，所以长期来看链把最多时间花在这些状态中是合理的。在上面的例子中，这说明
\[
  \mathbf s=\left(\frac4{14},\frac3{14},\frac2{14},\frac3{14},\frac2{14}\right)
\]
是该随机游走的平稳分布。

习题 20 会考察加权无向网络上的随机游走；每条边都有一个权重，从 $i$ 出发时，游走者按照与可用边权重成比例的概率选择去向。事实证明，这也是一条可逆马尔可夫链。更令人惊讶的是，每一条可逆马尔可夫链都可以表示为某个加权无向网络上的随机游走！$\square$

下面是无向网络上随机游走的一个具体例子。

\noindent\textbf{例 11.4.5（棋盘上的骑士）}\quad 考虑一个骑士在 $4\times4$ 棋盘上随机移动。

\begin{figure}[htbp]
  \centering
  \includegraphics[width=0.36\textwidth]{11-5-2.png}
\end{figure}

这 16 个格子按网格标记；例如，骑士当前位于 B3，左上角格子为 A4。骑士的每一步都是 L 形跳跃：先水平移动两格再垂直移动一格，或反过来。例如，从 B3 出发，骑士可以移动到 A1、C1、D2 或 D4；从 A4 出发，骑士可以移动到 B2 或 C3。注意，从浅色格子出发，骑士总会移动到深色格子，反之亦然。

假设每一步中，骑士等可能地随机选择一种可行走法。这给出了一条以 16 个格子为状态的马尔可夫链。求该链的平稳分布。

\noindent\textit{解.}\quad 棋盘上只有三类格子：4 个中心格、4 个角格（如 A4）以及 8 个边格（如 B4；这里不把角格算作边格）。可以把棋盘看作一个无向网络：若两个格子之间可以通过一次骑士移动到达，就在二者之间连一条边。于是中心格的度为 4，角格的度为 2，边格的度为 3，所以它们的平稳概率分别为 $4a,2a,3a$，其中 $a$ 为某个常数。

为了求 $a$，统计各类格子的数量，得到
\[
  4a\cdot4+2a\cdot4+3a\cdot8=1,
\]
从而 $a=1/48$。因此，每个中心格的平稳概率为 $4/48=1/12$，每个角格的平稳概率为 $2/48=1/24$，每个边格的平稳概率为 $3/48=1/16$。$\square$

第三，也是最后一类，如果一条马尔可夫链在每个时间段只能向左移动一步、向右移动一步，或留在原地，则称它为\textbf{生灭链}。所有生灭链都是可逆的。

\noindent\textbf{例 11.4.6（生灭链）}\quad 状态空间为 $\{1,2,\ldots,M\}$ 的\textbf{生灭链}是指转移矩阵 $Q=(q_{ij})$ 满足如下条件的马尔可夫链：若 $|i-j|=1$，则 $q_{ij}>0$；若 $|i-j|\ge2$，则 $q_{ij}=0$。这就是说，可以向左一步或向右一步移动（边界处除外），但不可能一步跳得更远。这个名称来自人口增长或减少的应用，其中向右一步可理解为出生，向左一步可理解为死亡。

例如，若下图中标出的转移均有正概率，则这条链就是生灭链；但从一个状态到自身的环允许有 0 概率。

\begin{figure}[htbp]
  \centering
  \includegraphics[width=0.68\textwidth]{11-5-3.png}
\end{figure}

现在证明任意生灭链都是可逆的，并构造其平稳分布。令 $s_1$ 为一个待定的正数。由于希望 $s_1q_{12}=s_2q_{21}$，令
\[
  s_2=s_1q_{12}/q_{21}.
\]
接着，由于希望 $s_2q_{23}=s_3q_{32}$，令
\[
  s_3=s_2q_{23}/q_{32}=s_1q_{12}q_{23}/(q_{32}q_{21}).
\]
继续这样做，令
\[
  s_j=\frac{s_1q_{12}q_{23}\cdots q_{j-1,j}}
  {q_{j,j-1}q_{j-1,j-2}\cdots q_{21}},
\]
对所有 $2\le j\le M$ 成立。选择 $s_1$ 使得所有 $s_j$ 之和为 1。于是该链相对于 $\mathbf s$ 可逆，因为当 $|i-j|\ge2$ 时 $q_{ij}=q_{ji}=0$，而当 $|i-j|=1$ 时，按构造有 $s_iq_{ij}=s_jq_{ji}$。因此，$\mathbf s$ 是平稳分布。$\square$

Ehrenfest 链是一条生灭链，可以用作气体分子扩散的简单模型。它的平稳分布会是一个二项分布。

\noindent\textbf{例 11.4.7（Ehrenfest）}\quad 有两个容器，总共有 $M$ 个可区分粒子。每次转移时，随机选取一个粒子，并把它从当前所在容器移到另一个容器。最初所有粒子都在第二个容器中。令 $X_n$ 为时刻 $n$ 第一个容器中的粒子数，因此 $X_0=0$，而从 $X_n$ 到 $X_{n+1}$ 的转移按上述方式进行。这是一条周期马尔可夫链，状态空间为 $\{0,1,\ldots,M\}$。

\begin{figure}[htbp]
  \centering
  \includegraphics[width=0.42\textwidth]{11-5-4.png}
\end{figure}

如上图所示，第一个容器有 $X_n$ 个粒子，第二个容器有 $M-X_n$ 个粒子。我们用可逆性条件说明，$\mathbf s=(s_0,s_1,\ldots,s_M)$，其中
\[
  s_i=\binom{M}{i}\left(\frac12\right)^M,
\]
是平稳分布。注意这正是 $\operatorname{Bin}(M,1/2)$ 的 PMF。

令 $s_i$ 如上，并检验 $s_iq_{ij}=s_jq_{ji}$。若 $j=i+1$（其中 $i<M$），则
\[
\begin{aligned}
  s_iq_{ij}
  &=
  \binom{M}{i}\left(\frac12\right)^M\frac{M-i}{M}
  =
  \frac{M!}{(M-i)!i!}\left(\frac12\right)^M\frac{M-i}{M}
  =
  \binom{M-1}{i}\left(\frac12\right)^M,\\
  s_jq_{ji}
  &=
  \binom{M}{j}\left(\frac12\right)^M\frac{j}{M}
  =
  \frac{M!}{(M-j)!j!}\left(\frac12\right)^M\frac{j}{M}
  =
  \binom{M-1}{j-1}\left(\frac12\right)^M
  =
  s_iq_{ij}.
\end{aligned}
\]
类似计算可知，若 $j=i-1$（其中 $i>0$），也有 $s_iq_{ij}=s_jq_{ji}$。对所有其他 $i,j$，有 $q_{ij}=q_{ji}=0$。因此，$\mathbf s$ 是平稳的。

二项分布是平稳分布的自然猜测，因为马尔可夫链运行很长时间之后，每个粒子大致同等可能在任一容器中。不过，该 PMF 并不会收敛到二项分布，因为链的周期为 2；当 $n$ 为偶数时，$X_n$ 必为偶数，当 $n$ 为奇数时，$X_n$ 必为奇数。

幸运的是，平稳分布的另一个解释在这里仍然成立：$s_i$ 是链长期来看花在状态 $i$ 中的时间比例。更精确地，令 $I_k$ 为链在时刻 $k$ 处于状态 $i$ 的指示变量，可以证明
\[
  \frac1n\sum_{k=0}^{n-1}I_k\to s_i
\]
当 $n\to\infty$ 时以概率 1 成立。$\square$

\section*{11.5\quad 总结}

马尔可夫链是满足马尔可夫性质的一列随机变量 $X_0,X_1,X_2,\ldots$；马尔可夫性质说的是：给定现在，过去与未来条件独立：
\[
  P(X_{n+1}=j\mid X_n=i,X_{n-1}=i_{n-1},\ldots,X_0=i_0)
  =
  P(X_{n+1}=j\mid X_n=i)
  =
  q_{ij}.
\]
转移矩阵 $Q=(q_{ij})$ 给出从任一状态一步移动到任一其他状态的概率。转移矩阵的第 $i$ 行是给定 $X_n=i$ 时 $X_{n+1}$ 的条件 PMF。转移矩阵的第 $n$ 次幂给出 $n$ 步转移概率。如果指定初始条件 $s_i=P(X_0=i)$ 并令 $\mathbf s=(s_1,\ldots,s_M)$，则 $X_n$ 的边缘 PMF 为 $\mathbf sQ^n$。

马尔可夫链的状态可以分为常返或暂留：若链会一再返回该状态，则为常返；若链最终会永远离开该状态，则为暂留。状态也可以按周期分类；状态 $i$ 的周期是从 $i$ 出发并返回 $i$ 所可能用到的步数的最大公约数。若能在有限步内从任一状态到达任一状态，则链是不可约的；若每个状态的周期都为 1，则链是非周期的。

有限马尔可夫链的平稳分布是一个满足 $\mathbf sQ=\mathbf s$ 的 PMF $\mathbf s$。在若干条件下，有限马尔可夫链的平稳分布存在且唯一，并且当 $n\to\infty$ 时 $X_n$ 的 PMF 收敛到 $\mathbf s$。若状态 $i$ 的平稳概率为 $s_i$，则从 $i$ 出发返回 $i$ 的期望时间为 $r_i=1/s_i$。

若一个 PMF $\mathbf s$ 满足可逆性条件
\[
  s_iq_{ij}=s_jq_{ji}
\]
对所有 $i,j$ 成立，则它保证 $\mathbf s$ 是转移矩阵 $Q=(q_{ij})$ 的马尔可夫链的平稳分布。存在满足可逆性条件的 $\mathbf s$ 的马尔可夫链称为可逆链。我们讨论了三类可逆链：

\begin{enumerate}
\item 如果转移矩阵是对称的，则平稳分布是在所有状态上的均匀分布。
\item 如果链是无向网络上的随机游走，则平稳分布与度序列成比例，即
\[
  s_j=\frac{d_j}{\sum_i d_i}.
\]
\item 如果链是生灭链，则平稳分布满足
\[
  s_j=\frac{s_1q_{12}q_{23}\cdots q_{j-1,j}}
  {q_{j,j-1}q_{j-1,j-2}\cdots q_{21}},
\]
其中 $j>1$，最后求出 $s_1$ 使 $s_1+\cdots+s_M=1$。
\end{enumerate}

\begin{figure}[htbp]
  \centering
  \includegraphics[width=0.88\textwidth]{11-6.png}
  \caption{给定转移矩阵 $Q$ 和状态上的分布 $\mathbf t$，可以通过按 $\mathbf t$ 选择 $X_0$，再按转移概率运行链，生成马尔可夫链 $X_0,X_1,\ldots$。一个重要事件是 $X_n=i$，即链在时刻 $n$ 访问状态 $i$。于是可以用 $Q$ 和 $\mathbf t$ 求出 $X_n$ 的 PMF，并且在本章讨论的条件下，该 PMF 会收敛到平稳分布 $\mathbf s$。若一开始就按 $\mathbf s$ 启动链，则链永远保持平稳。}
\end{figure}

图 11.6 比较了运行具有转移矩阵 $Q$ 的马尔可夫链的两种方式：按状态上的任意分布 $\mathbf t$ 选择初始状态，或按平稳分布 $\mathbf s$ 选择初始状态。在前一种情形中，经过 $n$ 步后的精确 PMF 可以用 $Q$ 和 $\mathbf t$ 表示，并且 PMF 会收敛到 $\mathbf s$（在本章讨论的一些很一般的条件下）。在后一种情形中，链永远保持平稳。

\section*{11.7\quad 习题}

标有 \textnormal{(S)} 的习题在 \texttt{http://stat110.net} 有详细解答。

\subsection*{马尔可夫性质}

\begin{enumerate}
\item[$1.$] \textnormal{(S)} 设 $X_0,X_1,X_2,\ldots$ 是一条马尔可夫链。证明 $X_0,X_2,X_4,X_6,\ldots$ 也是一条马尔可夫链，并直观解释为什么这件事合理。

\item[$2.$] \textnormal{(S)} 设 $X_0,X_1,X_2,\ldots$ 是一条不可约马尔可夫链，状态空间为 $\{1,2,\ldots,M\}$，其中 $M\ge3$，转移矩阵为 $Q=(q_{ij})$，平稳分布为 $\mathbf s=(s_1,\ldots,s_M)$。设初始状态 $X_0$ 服从平稳分布，即 $P(X_0=i)=s_i$。

(a) 平均来说，$X_0,X_1,\ldots,X_9$ 中有多少个等于 3？用 $\mathbf s$ 表示并化简。

(b) 令 $Y_n=(X_n-1)(X_n-2)$。当 $M=3$ 时，给出一个 $Q$ 的例子，使得 $Y_0,Y_1,\ldots$ 是马尔可夫链；再给出另一个 $Q$ 的例子，使得 $Y_0,Y_1,\ldots$ 不是马尔可夫链。在例子中，要求所有 $q_{ij}>0$，并保证最终可以从任一状态到达任一状态。

\item[$3.$] 一条马尔可夫链有两个状态 $A$ 与 $B$，转移如下：

\begin{figure}[htbp]
  \centering
  \includegraphics[width=0.38\textwidth]{11-6-1.png}
\end{figure}

假设我们不能直接观察这条马尔可夫链，称其为 $X_0,X_1,X_2,\ldots$。相反，每当链从 $A$ 回到 $A$ 时，我们观察到 0；每当它改变状态时，我们观察到 1。把由 0 与 1 组成的序列称为 $Y_0,Y_1,Y_2,\ldots$。

例如，如果 $X$ 链开始为
\[
  A,A,B,A,B,A,A,\ldots,
\]
则 $Y$ 链开始为
\[
  0,1,1,1,1,0,\ldots.
\]

(a) 证明 $Y_0,Y_1,Y_2,\ldots$ 不是马尔可夫链。

(b) 在例 11.1.3 中，我们通过扩大状态空间以纳入二阶依赖，处理了马尔可夫性质被破坏的情形。证明这样的技巧对 $Y_0,Y_1,Y_2,\ldots$ 不起作用。也就是说，无论 $m$ 多大，
\[
  Z_n=(Y_{n-m+1},Y_{n-m+2},\ldots,Y_n),\qquad n=m-1,m,\ldots
\]
仍然不是马尔可夫链。

\item[$4.$] 有三块漂浮在熔岩海上的石块，从左到右标为 1、2、3。电子游戏角色 Mark the Kangaroo 站在石块 1 上。为了到达安全地带，他必须到达石块 3。他不能直接从石块 1 跳到石块 3；唯一办法是先从石块 1 跳到石块 2，再从石块 2 跳到石块 3。每次跳跃时，他成功的概率为 $1/2$，以 $1/2$ 的概率掉进熔岩而“死亡”。如果他“死亡”，就从石块 1 重新开始。

令 $J$ 为 Mark 为了到达石块 3 所做跳跃的总次数。

(a) 用马尔可夫性质与条件期望求 $E(J)$。

(b) 说明这个问题如何与第 9 章中的一个抛硬币问题相关。

\subsection*{平稳分布}

\item[$5.$] \textnormal{(S)} 考虑下图所示的马尔可夫链，其中 $0<p<1$，箭头上的标号表示转移概率。

\begin{figure}[htbp]
  \centering
  \includegraphics[width=0.42\textwidth]{11-6-2.png}
\end{figure}

(a) 写出该链的转移矩阵 $Q$。

(b) 求该链的平稳分布。

(c) 当 $n\to\infty$ 时，$Q^n$ 的极限是什么？

\item[$6.$] \textnormal{(S)} 考虑下图所示马尔可夫链，状态空间为 $\{1,2,3,4\}$，箭头右侧的标号表示转移概率。

\begin{figure}[htbp]
  \centering
  \includegraphics[width=0.48\textwidth]{11-6-3.png}
\end{figure}

(a) 写出该链的转移矩阵 $Q$。

(b) 哪些状态（如果有）是常返的？哪些状态（如果有）是暂留的？

(c) 求该链的两个不同平稳分布。

\item[$7.$] \textnormal{(S)} Daenerys 有三条龙：Drogon、Rhaegal 与 Viserion。每条龙独立地在世界中探索，寻找美味食物。令 $X_n,Y_n,Z_n$ 分别为时刻 $n$ Drogon、Rhaegal、Viserion 的位置；时间是离散的，可能位置的数量是有限数 $M$。它们的路径 $X_0,X_1,X_2,\ldots$，$Y_0,Y_1,Y_2,\ldots$ 和 $Z_0,Z_1,Z_2,\ldots$ 是相互独立的马尔可夫链，具有相同的平稳分布 $\mathbf s$。每条龙的初始位置都按平稳分布随机生成。

(a) 令状态 0 为家（所以 $s_0$ 是家的平稳概率）。求直到时刻 24 为止 Drogon 在家的次数的期望，即 $X_0,X_1,\ldots,X_{24}$ 中有多少个为状态 0 的期望，用 $s_0$ 表示。

(b) 若要同时追踪 3 条龙，需要考虑位置向量 $(X_n,Y_n,Z_n)$。该向量有 $M^3$ 个可能取值；假设每个取值都被编号为 $1$ 到 $M^3$。例如，当 $M=2$ 时，可以把状态 $(0,0,0),(0,0,1),(0,1,0),\ldots,(1,1,1)$ 分别编码为 $1,2,3,\ldots,8$。令 $W_n$ 为表示 $(X_n,Y_n,Z_n)$ 的 $1$ 到 $M^3$ 之间的数字。判断 $W_0,W_1,\ldots$ 是否为马尔可夫链。

(c) 已知三条龙在时刻 0 都从家出发，求三条龙再次同时在家的期望时间。

\item[$8.$] 考虑如下马尔可夫链，它有 $52!\approx8\times10^{67}$ 个状态。状态是一副标准 52 张牌的所有可能排列。链每走一步，就从牌堆中选取两张不同的牌（所有牌对等可能），并交换这两张牌。求该链的平稳分布。

\subsection*{可逆性}

\item[$9.$] \textnormal{(S)} 一条状态空间为 $\{-3,-2,-1,0,1,2,3\}$ 的马尔可夫链 $X_0,X_1,\ldots$ 如下运行。链从 $X_0=0$ 开始。若 $X_n$ 不是端点（$-3$ 或 $3$），则 $X_{n+1}$ 以各 $1/2$ 的概率为 $X_n-1$ 或 $X_n+1$。否则链会从端点反射：从 3 必到 2，从 $-3$ 必到 $-2$。该链如下图所示。

\begin{figure}[htbp]
  \centering
  \includegraphics[width=0.58\textwidth]{11-6-4.png}
\end{figure}

(a) $|X_0|,|X_1|,|X_2|,\ldots$ 是否也是马尔可夫链？解释。

提示：对 (a) 和 (b)，思考给定现在时过去与未来是否条件独立；不要对 $7\times7$ 转移矩阵做计算！

(b) 令 $\operatorname{sgn}$ 为符号函数：若 $x>0$，则 $\operatorname{sgn}(x)=1$；若 $x<0$，则 $\operatorname{sgn}(x)=-1$；且 $\operatorname{sgn}(0)=0$。$\operatorname{sgn}(X_0),\operatorname{sgn}(X_1),\operatorname{sgn}(X_2),\ldots$ 是否为马尔可夫链？解释。

(c) 求链 $X_0,X_1,X_2,\ldots$ 的平稳分布。

(d) 找一种简单方法，修改某些 $i\in\{-3,3\}$ 对应的转移概率 $q_{ij}$，使修改后链的平稳分布在所有状态上均匀。

\item[$10.$] \textnormal{(S)} 令 $G$ 为一个无向网络，其节点标为 $1,2,\ldots,M$（不允许自环），其中 $M\ge2$，且该网络上的随机游走不可约。令 $d_j$ 为节点 $j$ 的度。构造状态空间为 $1,2,\ldots,M$ 的马尔可夫链如下：从状态 $i$ 出发，先在所有与 $i$ 有边相连的节点中均匀随机选一个候选节点 $j$；然后以概率 $\min(d_i/d_j,1)$ 到达 $j$，否则留在 $i$。

(a) 求该链从 $i$ 到 $j$ 的转移概率 $q_{ij}$，对所有状态 $i,j$ 都给出。

(b) 求该链的平稳分布。

\item[$11.$] \textnormal{(S)} (a) 考虑状态空间为 $\{1,2,\ldots,7\}$ 的马尔可夫链，状态按下图的“圆圈”排列，转移由顺时针或逆时针移动一步给出，二者等概率。例如，从状态 6 出发，链以各 $1/2$ 的概率到状态 7 或状态 5；从状态 7 出发，链以各 $1/2$ 的概率到状态 1 或状态 6。链从状态 1 开始。

\begin{figure}[htbp]
  \centering
  \includegraphics[width=0.35\textwidth]{11-6-5.png}
\end{figure}

求该链的平稳分布。

(b) 考虑一条通过“展开圆圈”得到的新链。现在状态按下图排列。从状态 1 出发，链总是到状态 2；从状态 7 出发，链总是到状态 6。求新的平稳分布。

\begin{figure}[htbp]
  \centering
  \includegraphics[width=0.65\textwidth]{11-6-6.png}
\end{figure}

\item[$12.$] \textnormal{(S)} 令 $X_n$ 为某只股票第 $n$ 天开始时的价格，并假设 $X_0,X_1,X_2,\ldots$ 是转移矩阵为 $Q$ 的马尔可夫链。（为简单起见，假设股票价格永远不会低于 0，也不会超过某个上界，并且总是四舍五入到最接近的美元。）

(a) 一个懒惰投资者每年只查看一次股票，观察第 $0,365,2\cdot365,3\cdot365,\ldots$ 天的价格。因此投资者观察到 $Y_0,Y_1,\ldots$，其中 $Y_n$ 是 $n$ 年后（即 $365n$ 天后）的价格。$Y_0,Y_1,\ldots$ 是否也是马尔可夫链？解释原因；若是，它的转移矩阵是什么？

(b) 股票价格总是在 $0$ 美元到 $28$ 美元之间的整数。每天到下一天，股价会上涨或下跌 1 美元或 2 美元，这些可能性等概率（但当股票处于或接近边界，即 $0$、$1$、$27$ 或 $28$ 美元时除外）。

如果股票在 0 美元，下一天会上涨到 1 美元或 2 美元（得到政府救助之后）。如果股票在 28 美元，下一天会跌到 27 美元或 26 美元。如果股票在 1 美元，它要么上涨到 2 或 3 美元，要么下跌到 0 美元（等概率）；类似地，如果股票在 27 美元，它要么上涨到 28 美元，要么下跌到 26 或 25 美元。求该链的平稳分布。

\item[$13.$] \textnormal{(S)} 在国际象棋中，王可以每次沿任意方向（水平、竖直或对角）移动一格。

\begin{figure}[htbp]
  \centering
  \includegraphics[width=0.36\textwidth]{11-6-7.png}
\end{figure}

例如，在图中，从当前位置出发，王可以移动到 8 个可能格子中的任一个。一个王在一块其余为空的 $8\times8$ 棋盘上游走，每一步在所有可能走法中等可能选择。求该链的平稳分布（当然，不要显式列出长度为 64 的向量！把 64 个格子分类，并说明每类格子的平稳概率）。

\item[$14.$] 一个棋子在一块其余为空的 $8\times8$ 棋盘上游走。每一步中，该棋子（王、后、车、象或马）在合法走法中均匀随机选择去向（按国际象棋规则；若不熟悉这些规则，需要自行查阅）。

(a) 对这些情形分别判断该马尔可夫链是否不可约，以及是否非周期。

马的提示：注意马的每一步总是从浅色格子到深色格子，或反过来。骑士巡游是棋盘上的一列马步，使得马恰好访问每个格子一次。存在许多骑士巡游。

(b) 本小问假设棋子是车，初始位置均匀随机选择。求 $n$ 步之后车所在位置的分布。

(c) 现在假设棋子是王，初始位置确定为左上角。求王回到该角落所需步数的期望；答案应充分化简，最好不超过 140 个字符。

(d) 上一小问中王的随机游走的平稳分布不是棋盘 64 个格子上的均匀分布。下面给出一种修改链的方法，以得到均匀平稳分布。把格子标为 $1,2,\ldots,64$，令 $d_i$ 为从格子 $i$ 出发的合法走法数量。假设王当前在格子 $i$。链的下一步如下决定：

第 1 步：在从 $i$ 出发的合法走法中均匀随机选一个候选格子 $j$。

第 2 步：以概率 $\min(d_i/d_j,1)$ 掷出正面。若为正面，则到 $j$；否则留在 $i$。

证明修改后的链在 64 个格子上的平稳分布是均匀的。

\item[$15.$] \textnormal{(S)} 不使用矩阵，求下图所示马尔可夫链的平稳分布。箭头上方的数字为相应转移概率。

\begin{figure}[htbp]
  \centering
  \includegraphics[width=0.72\textwidth]{11-6-8.png}
\end{figure}

\item[$16.$] 有两个盒子，总共有 $2N$ 个可区分球。最初，第一个盒子有 $N$ 个白球，第二个盒子有 $N$ 个黑球。每一阶段，从每个盒子中随机取一个球并交换。令 $X_n$ 为时刻 $n$ 第一个盒子中的黑球数。这是状态空间 $\{0,1,\ldots,N\}$ 上的马尔可夫链。

(a) 给出该链的转移概率。

(b) 通过验证可逆性条件，证明 $(s_0,s_1,\ldots,s_N)$ 是平稳分布，其中
\[
  s_i=\frac{\binom Ni\binom N{N-i}}{\binom{2N}{N}}.
\]

\item[$17.$] 求状态空间为 $\{0,1,\ldots,110\}$ 的马尔可夫链 $X_0,X_1,X_2,\ldots$ 的平稳分布，其转移概率由
\[
  P(X_{n+1}=j\mid X_n=0)=p,\qquad j=1,2,\ldots,110,
\]
\[
  P(X_{n+1}=0\mid X_n=0)=1-110p,
\]
\[
  P(X_{n+1}=j\mid X_n=j)=1-r,\qquad j=1,2,\ldots,110,
\]
\[
  P(X_{n+1}=0\mid X_n=j)=r,\qquad j=1,2,\ldots,110,
\]
给出，其中 $p$ 与 $r$ 是常数，满足 $0<p<1/110$ 且 $0<r<1$。

\item[$18.$] 判断下图所示马尔可夫链是否可逆，并求该链的平稳分布。箭头左侧的标号给出相应转移概率。

\begin{figure}[htbp]
  \centering
  \includegraphics[width=0.58\textwidth]{11-6-9.png}
\end{figure}

\item[$19.$] Nausicaa Distribution 在 Etsy 上销售分布毛绒玩具。他们有 Evil Cauchy 毛绒玩具的两张不同照片，但不知道哪张更能促使顾客购买。每位访问网站的人会看到两张照片之一（称为照片 A 与照片 B），然后该访问者要么购买 Evil Cauchy（“成功”），要么不购买（“失败”）。

令 $a$ 与 $b$ 分别为展示照片 A 与照片 B 时成功的概率。即使 Evil Cauchy 无法抗拒，也假设 $0<a<1$ 且 $0<b<1$。采用如下策略（注意该策略不需要知道 $a$ 与 $b$）。第一个访问者看照片 A。如果该访问者购买了 Evil Cauchy，则下一位访问者继续看照片 A；否则切换到照片 B。类似地，若第 $n$ 位访问者是“成功”，则第 $n+1$ 位访问者看同一张照片，否则切换到另一张照片。

(a) 说明如何把所得过程表示为马尔可夫链，画出图并给出转移矩阵。状态为 A1、B1、A0、B0（转移矩阵和平稳分布都按这个顺序），其中例如状态 A1 表示当前访问者看的是照片 A 且成功。

(b) 判断该链是否可逆。

提示：先思考哪些转移概率为 0，哪些不为 0。

(c) 证明平稳分布与
\[
  \left(\frac{a}{1-a},\frac{b}{1-b},1,1\right)
\]
成比例，并求出平稳分布。

(d) 证明当 $a\ne b$ 时，每位访问者成功的平稳概率严格好于每次独立、等概率随机选择展示哪张照片所得到的成功概率。

\item[$20.$] 本题考察加权无向网络上的随机游走。假设给定一个无向网络，其中每条边 $(i,j)$ 都有一个非负权重 $w_{ij}$（允许 $i=j$）。假设 $w_{ij}=w_{ji}$，因为从 $i$ 到 $j$ 的边与从 $j$ 到 $i$ 的边被视为同一条边。为简化记号，若 $(i,j)$ 不是一条边，则定义 $w_{ij}=0$。

当处于节点 $i$ 时，下一步按与权重成比例的概率选择连接到 $i$ 的一条边。例如，如果游走处于节点 1，且从节点 1 出发有 3 条可能的边，权重分别为 7、1、4，则第一条边以概率 $7/12$ 被走过，第二条以概率 $1/12$ 被走过，第三条以概率 $4/12$ 被走过。若所有权重都等于 1，则该过程退化为我们先前研究的网络随机游走。

(a) 令 $v_i=\sum_j w_{ij}$，对所有节点 $i$ 成立。证明节点 $i$ 的平稳分布与 $v_i$ 成比例。

(b) 证明每一条可逆马尔可夫链都可以表示为某个加权无向网络上的随机游走。也就是说，给定一条可逆马尔可夫链的转移矩阵 $Q$，证明可以选择权重 $w_{ij}$，使得上面定义的随机游走是一条转移矩阵为 $Q$ 的马尔可夫链。务必检查 $w_{ij}=w_{ji}$。

提示：令 $w_{ij}=s_iq_{ij}$，其中 $\mathbf s$ 是平稳分布，$q_{ij}$ 是 $Q$ 的 $(i,j)$ 元。

\subsection*{综合练习}

\item[$21.$] \textnormal{(S)} 一只猫和一只老鼠独立地在两个房间之间来回移动。每个时间步中，猫以概率 0.8 从当前房间移动到另一个房间。若从房间 1 出发，老鼠以概率 0.3 移动到房间 2（否则留在原处）；若从房间 2 出发，老鼠以概率 0.6 移动到房间 1（否则留在原处）。

(a) 求猫链与老鼠链的平稳分布。

(b) 注意共有 4 个可能的（猫，老鼠）状态：二者都在房间 1、猫在房间 1 且老鼠在房间 2、猫在房间 2 且老鼠在房间 1、二者都在房间 2。分别把这些情形编号为 1、2、3、4，并令 $Z_n$ 为表示时刻 $n$ 的（猫，老鼠）状态的编号。$Z_0,Z_1,Z_2,\ldots$ 是否为马尔可夫链？

(c) 现在假设若猫和老鼠在同一房间，猫就会吃掉老鼠。我们想知道两种初始配置下猫吃掉老鼠之前的期望时间（步数）：猫从房间 1、老鼠从房间 2 出发，以及反过来。建立一个含两个未知数的二元线性方程组，其解就是所需值。

\item[$22.$] (a) Alice 与 Bob 相互独立地在一所含有 $M$ 个房间的房子中随机游走，房间标为 $1,2,\ldots,M$。令 $d_i$ 为房间 $i$ 的门数（通往其他房间，不通向室外）。每一步中，Alice 等可能地随机选择一扇门穿过，Bob 也独立地这样做。他们各自遵循的马尔可夫链不可约且非周期。令 $A_n$ 与 $B_n$ 分别为时刻 $n$ Alice 与 Bob 所在房间，$n=0,1,2,\ldots$。求
\[
  \lim_{n\to\infty}P(A_n=i,B_n=j).
\]

(b) 沿用 (a) 的设定，令 $p_{ij}$ 为从房间 $i$ 到房间 $j$ 的转移概率。令 $t_{ik}$ 为 Alice 与 Bob 第一次同处一个房间的期望时间，其中 Alice 从房间 $i$ 出发，Bob 从房间 $k$ 出发（注意若 $i=k$，则 $t_{ik}=0$）。

给出一个线性方程组，其解会给出所有房间 $i,k$ 对应的 $t_{ik}$。

提示：对 Alice 与 Bob 的第一次移动条件化。

\item[$23.$] 令 $\{X_n\}$ 是状态空间 $\{0,1,2\}$ 上的马尔可夫链，转移矩阵为
\[
  Q=
  \begin{pmatrix}
    0.8 & 0.2 & 0\\
    0 & 0.8 & 0.2\\
    0 & 0 & 1
  \end{pmatrix}.
\]
链从 $X_0=0$ 开始。令 $T$ 为到达状态 2 所需的时间：
\[
  T=\min\{n:X_n=2\}.
\]
通过画出马尔可夫链并讲述一个故事，求 $E(T)$ 与 $\operatorname{Var}(T)$。

\item[$24.$] 考虑状态空间为 $\{1,2,3,4,5,6\}$ 的如下马尔可夫链。

\begin{figure}[htbp]
  \centering
  \includegraphics[width=0.48\textwidth]{11-6-10.png}
\end{figure}

(a) 假设链从状态 1 开始。求该链返回状态 1 的次数的分布。

(b) 长期来看，链把多少比例的时间花在状态 3？简要解释。

\item[$25.$] 令 $Q$ 为状态空间 $\{1,2,\ldots,M\}$ 上某马尔可夫链的转移矩阵，其中状态 $M$ 是一个\textbf{吸收状态}，也就是说，从状态 $M$ 出发，链永远不会离开。假设从任一其他状态出发，都可能在若干步内到达 $M$。

(a) 哪些状态是常返的，哪些状态是暂留的？解释。

(b) 当 $n\to\infty$ 时，$Q^n$ 的极限是什么？

(c) 对 $i,j\in\{1,2,\ldots,M-1\}$，已知链在时刻 0 处于状态 $i$，求链在时刻 $n$ 处于状态 $j$ 的概率（答案应用 $Q$ 表示）。

(d) 对 $i,j\in\{1,2,\ldots,M-1\}$，已知链在时刻 0 处于状态 $i$，求直到并包括时刻 $n$ 为止链处于状态 $j$ 的期望次数（用 $Q$ 表示）。

(e) 令 $R$ 为从 $Q$ 中删除最后一行与最后一列得到的 $(M-1)\times(M-1)$ 矩阵。证明 $(I-R)^{-1}$ 的 $(i,j)$ 元等于链从状态 $i$ 出发时，在吸收之前处于状态 $j$ 的期望次数。

提示：有
\[
  I+R+R^2+\cdots=(I-R)^{-1},
\]
类似于几何级数。此外，如果把 $Q$ 分块为
\[
  Q=
  \begin{pmatrix}
    R & B\\
    0 & 1
  \end{pmatrix},
\]
其中 $B$ 是 $(M-1)\times1$ 矩阵，$0$ 是 $1\times(M-1)$ 零矩阵，则
\[
  Q^k=
  \begin{pmatrix}
    R^k & B_k\\
    0 & 1
  \end{pmatrix}
\]
其中 $B_k$ 是某个 $(M-1)\times1$ 矩阵。

\item[$26.$] 在称为 Chutes and Ladders 的游戏中，玩家努力成为第一个到达棋盘上某个目的地的人。棋盘是由方格组成的网格，方格编号从 1 到总方格数。棋盘上有一些“滑梯”和“梯子”，每一个都连接一对方格。这里考虑单人版本（这可以直接扩展到多人版本，因为若有多名玩家，他们只是独立轮流行动，直到有人到达目的地）。

每一轮中，玩家掷一个公平骰子，骰子决定在网格上向前移动多少格；例如，若玩家在方格 5 且掷出 3，则前进到方格 8。若得到的方格是梯子的底端，则玩家爬上梯子，立刻到达更靠前的方格。若得到的方格是滑梯的顶端，则玩家立刻滑到滑梯底端。

这个游戏可以自然地看作马尔可夫链：给定玩家当前在棋盘上的位置，过去历史对于计算例如未来 3 步内获胜的概率并不重要。

考虑一个简化版 Chutes and Ladders，在下图所示的 $3\times3$ 棋盘上进行。玩家从方格 1 开始，目标是到达方格 9。每次移动时，掷一枚公平硬币；若为正面，玩家前进 1 格；若为反面，玩家前进 2 格。不过，棋盘上有 2 条梯子（上箭头）和 2 条滑梯（下箭头）。

\begin{figure}[htbp]
  \centering
  \includegraphics[width=0.36\textwidth]{11-6-11.png}
\end{figure}

(a) 解释为什么尽管有 9 个方格，我们可以用下面的 $5\times5$ 转移矩阵表示该游戏：
\[
  Q=
  \begin{pmatrix}
    0 & 0.5 & 0.5 & 0 & 0\\
    0 & 0 & 0.5 & 0.5 & 0\\
    0.5 & 0 & 0 & 0.5 & 0\\
    0.5 & 0 & 0 & 0 & 0.5\\
    0 & 0 & 0 & 0 & 1
  \end{pmatrix}.
\]

(b) 不使用矩阵或繁琐计算，求玩家访问方格 7 的次数的均值与方差。

(c) 求游戏持续时间的中位数（其中持续时间定义为掷硬币次数）。

提示：把持续时间的 CDF 与 $Q$ 的幂联系起来。

(d) 求游戏持续时间的均值（持续时间定义同上）。

提示：把持续时间与在暂留状态中花费的总时间联系起来，并使用上一题 (e)。
\end{enumerate}

\end{document}
